\documentclass[12pt, a4paper]{article}
\usepackage[T1]{fontenc}
\renewcommand{\rmdefault}{phv} % Arial
\renewcommand{\sfdefault}{phv} % Arial
% \usepackage{fontspec}
% \usepackage{uarial}
% \setmainfont{Times New Roman}
\usepackage[utf8]{inputenc}
\usepackage[
    a4paper,
    left=15mm,
    right=15mm,
    top=18mm,
    bottom=10mm
]{geometry}
\usepackage[
    skip=5pt,
    indent=0pt,
]{parskip}
\setlength{\headheight}{25pt}
\usepackage{setspace}
\fontsize{12pt}{24pt}\selectfont
\linespread{2} 

\usepackage{natbib}
\usepackage{hyperref}
\usepackage{xcolor}
\usepackage{xspace}
\usepackage{fancyhdr}
\usepackage[flushmargin, bottom]{footmisc}
\addtolength{\footnotesep}{-10mm}
\hypersetup{
    colorlinks,
    linkcolor={red!50!black},
    citecolor={blue!50!black},
    urlcolor={blue!80!black}
}

\newcommand{\HRule}{\rule{\linewidth}{0.5mm}}
\newcommand{\Hrule}{\rule{\linewidth}{0.3mm}}

% Project specific macros
\newcommand{\MFG}{Multi-Period Compliance MFG with Deep FBSDE Solvers\xspace}
\newcommand{\Report}{\href{https://github.com/OrangeAoo/Multi-Compliance-MFG-FBSDEs/blob/ca6adb57636c1fc1ff9daea8d57f004e5762fd4d/FinalReports/Report-v3.pdf}{\textbf{the Report}}\xspace}
\newcommand{\GitHubMFG}{\href{https://github.com/OrangeAoo/Multi-Period-Compliance-MFG-FBSDEs}{\textbf{my GitHub repository}}\xspace}
\newcommand{\Lowcarb}{Low-carbon Transition of China's Power System\xspace} 
\newcommand{\InsurP}{Insurance Pricing And Risk Modeling with ML\xspace}

% Intern specific macros
\newcommand{\CW}{Chongwen Quantitative\footnote{Beijing, China. Worked remotely.}\xspace}
\newcommand{\TQ}{TrexQuant\footnote{Stamford, CT, United States. Worked remotely.}\xspace}
\newcommand{\Higgs}{Higgs Asset\footnote{Hangzhou, China. Worked onsite.}\xspace}

% School specific macros
\newcommand{\NJU}{Nanjing University\xspace}
\newcommand{\THU}{Tsinghua University\xspace}
\newcommand{\TCJC}{Tsinghua-CTG Joint Center\xspace}
\newcommand{\CU}{Columbia University\xspace}
\newcommand{\schoolShort}{NYU Tandon\xspace}
\newcommand{\schoolLong}{NYU Tandon School of Engineering\xspace}
\newcommand{\programShort}{FRE-ECE\xspace}
\newcommand{\programLong}{Department of Finance and Risk Engineering - Electrical Engineering\xspace}

% Profs specific macros
\newcommand{\SC}{\href{https://www.stevenacampbell.com/}{Prof. Steven Campbell}\xspace}
\newcommand{\NT}{\href{https://engineering.nyu.edu/faculty/nizar-touzi}{Prof. Nizar Touzi}\xspace}
\newcommand{\RX}{\href{https://engineering.nyu.edu/faculty/renyuan-xu}{Prof. Renyuan Xu}\xspace}
\newcommand{\AT}{Prof. Agnes Tourin\xspace}
\newcommand{\AAb}{Prof. Amine Mohamed Aboussalah\xspace}
\newcommand{\YQ}{Yang Qu (Postdoctoral Fellow)\xspace}
\newcommand{\CA}{\href{https://www.linkedin.com/in/carocha/?originalSubdomain=ch}{Carlos Arocha}\xspace}
\newcommand{\LX}{\href{https://sps.columbia.edu/person/lina-xu}{Prof. Lina Xu}}
\newcommand{\YW}{\href{https://sps.columbia.edu/person/yubo-wang-phd}{Prof. Yubo Wang}\xspace}
\newcommand{\Tom}{\href{https://sps.columbia.edu/person/thomas-murphy}{Prof. Thomas Murphy}\xspace}
\newcommand{\Ira}{\href{https://sps.columbia.edu/person/ira-kastrinsky}{Prof. Ira Kastrinsky}\xspace}

% Creates header for each page
\usepackage{fancyhdr}
\pagestyle{fancy}
\fancyhf{}
\fancyhead[R]{\header\hskip\linepagesep\vfootline\thepage}
\newskip\linepagesep \linepagesep 5pt\relax
\def\vfootline{%
    \begingroup
    	\rule[-10pt]{0.75pt}{25pt}
    \endgroup
}
\def\header{%
\small
\begin{minipage}{120pt}
    \hfill Orange Ao 	% Applicant Name
   \par \hfill FRE, PhD, Fall 2025 % Area, Program, Cycle, Year
\end{minipage}
}
\fancyhead[L]{Personal Essay | \schoolShort}
\renewcommand\headrulewidth{0.2pt}

\begin{document}

I am a marathoner. “Why the hell am I doing this?” was the question I’ve asked myself a thousand times on the track. And I know it will be asked a thousand more times during my PhD. As an undergrad from outside the States, emphasizing achievements like publications might not serve me well compared to other applicants. So rather than highlight my current \textit{ordinate}, I'd rather focus on the \textit{first and second-order derivatives} of my curve and where it leads to in this essay.

A year ago, I made a bold decision to independently apply for an exchange at \CU at the cost of suspending my studies at \NJU (NJU), as it didn't officially cooperate with Columbia. I also risked not earning enough credits to graduate, since they were not eligible for transfer back to NJU\footnote{This would explain the gap on my NJU transcript in Spring 2024.}. Nonetheless, I pursued this path with the belief that the greater resources at Columbia would make it cost-beneficial. Aware of the heightened risks compared to typical exchange students, I was more driven to excel in courses and to seize every opportunity that came my way. And I was right - had I not spent the 8 months at Columbia, I'd never have made the invaluable connections with the professors who helped uncover my potential for PhD research.

So I want to highlight my research processes since then. Exposed to an unfamiliar field of insurance pricing, I picked up the basics quickly and incorporated my strength in ML into the project. Mentored by actuary \CA\footnote{Also under the supervision of \LX, \YW, \Tom, and \Ira.}, I adopted K-Means to identify risk clusters and customized exposure rating models for each cluster - an innovative pricing approach that has yet to be implemented in the industry. This further inspired me to develop a Peer Grouping System as a quantitative researcher intern, combining the Affinity Propagation and XGBoost to provide a data-driven alternative for the industry grouping that is better suited for high-frequency strategies. 

\SC introduced me to stochastic optimization. Under his supervision, I took an independent approach to mastering PyTorch and developed Deep FBSDE Solvers into user-friendly packages\footnote{The results are summarized in \Report and made open-source on \GitHubMFG.} in 2 months. I especially loved the philosophical lessons from modeling how defaultable agents with different strategies responded to penalty policies: don't be blinded by the current \textit{ordinate}, always plan ahead, and invest in the \textit{second-order derivate} of growth early. 

In my analogy, a PhD degree is just another \textit{ordinate} representing my track record and new starting point. The \textit{first-order derivative} is the resources such as knowledge and connections gathered during my PhD, indicating the speed I achieve my future goal. And the \textit{second-order derivative} determines the rate of expanding resources - traits like simultaneous persistence and flexibility honed through years of toil and grind, the capability to marshal resources trained through fundraising, and the alertness to spot and seize opportunities.

True, a PhD entails long years, low pay, and often, frustrations of no immediate progress. “So why?” As a marathoner focused on the finish line, I’m after something big - building an avant-garde business. And I am more afraid of the prescribed path with sublinear growth as I age. While I can’t predict exactly what steps I will take as a “5-year plan” since there is randomness in everything and the nature of my business may evolve based on market conditions, I do anticipate almost surely the exponential curve with a good starting point upon completing my PhD.

% Simply put, it’s because with greater risk and effort comes greater reward - and I am after something big. 

As the only program with industry professors to my knowledge, Tandon FRE stands at the top of my list for its strong industry connections, signaling a great \textit{first-order derivative}, i.e. resources I can gain. Also, this program strikes a perfect balance between breadth and depth: the faculty's expertise spans a wide range of topics, yet each professor focuses deeply on a specific area. I am particularly drawn to \NT’s theoretical study of the Principle-Agent problem, as well as \RX's extensive research on generative models applied to stochastic controls. These areas not only align closely with my research interests but also will equip me with the fundamental tools to model the world of chaos.

In summary, my PhD is not just about earning a degree, but about acquiring the resources and traits essential to build something transformative. With a strength in stochastic modeling and deep connections to industry, I will be well-positioned to achieve my entrepreneurial goal in the dynamic markets. I look forward to the upcoming marathon at Tandon FRE.

\end{document}

% That's All Folks.

% Best of luck, you got this! :)